% Volume V: Admissibility and Continuation
% Part X. Public Institutions as Epistemic Ecologies
% Chapter 55: Science and Adjudication
% Spine: proof-spines/ch055-spine.md

\chapter{Science and Adjudication}
\label{ch:ch055}

Scientific conclusions remain revisable, while legal decisions often require finite closure. The chapter therefore distinguishes replication, peer review, evidentiary hearings, precedent, and final judgment rather than treating science and adjudication as interchangeable truth machines.

Its central question is when scientific uncertainty is imported badly into law and when legal demands for binary resolution distort scientific evidence. Science can leave a matter open while accumulating better models; adjudication must often decide a case, allocate responsibility, and stop. The problem is not whether one institution should dominate the other, but how revisable inquiry is translated into deadline-bound public judgment.
